\paragraph{Two further cache-flexible source variants.}
The derivation compiler can allow the numerical cache $D$ to have an
arbitrary value after a successful MP: each compiled premise in the
MP clause is followed by PUSH, and the generalization word also begins with PUSH.
Those words clear the incoming caches.  Two concrete source variants
were therefore tested.  The first tests $b=1$ with two decrements and
retains $D=\max(b-2,0)$ on a nonhalting success.  The second makes the
consequent $b=d+x+1$, so a zero raw cache marks the contradictory
consequent, and retains $D=b-2$ on a nonhalting success.  Each has
3,782 primitive sites but 163 distinct internal nodes, one more than
the initial halt-at-MP source and four more than PA196.  The complete
MP fragments pass 2,286 and 2,304 bounded exact-integer tests,
respectively, with all nine registers compared.  These are source
experiments; neither yields a smaller physical transition table, and
the finite tests are not proof-assistant certificates.  The supplied
\texttt{cache\_flexible\_variants.py} reconstructs both from the pinned
PA196 source and reproduces the receipts.
